\documentclass[12pt,a4paper]{article}
\usepackage[utf8]{inputenc}
\usepackage[indonesian]{babel}
\usepackage{amsmath}
\usepackage{amsfonts}
\usepackage{amssymb}
\usepackage{geometry}

\geometry{margin=2.5cm}

\title{Penyelesaian Titik Berat Benda Putar}
\author{Ahmad Fakhrul Bawani}
\date{}

\begin{document}

\maketitle

\section*{Soal}
Diketahui sebuah kurva dengan persamaan $x = y^{1/3}$ pada batas $y \in [1, 8]$ diputar terhadap \textbf{sumbu-x}. 

Karena benda diputar terhadap sumbu-x, maka berdasarkan prinsip simetri, koordinat titik berat pada sumbu-y adalah:
\[\bar{y} = 0\]

Untuk mencari koordinat $\bar{x}$, kita gunakan representasi fungsi $y$ terhadap $x$ agar proses integrasi lebih mudah diselesaikan:
\begin{itemize}
  \item Kurva: $y = x^3$
  \item Batas integral terhadap $x$: Jika $y \in [1, 8]$, maka $x \in [1, 2]$
  \item Turunan fungsi: $\frac{dy}{dx} = 3x^2$
  \item Elemen panjang busur ($ds$): 
  \[ds = \sqrt{1 + \left(\frac{dy}{dx}\right)^2} \, dx = \sqrt{1 + 9x^4} \, dx\]
\end{itemize}

---

\subsection*{1. Rumus Utama Titik Berat ($\bar{x}$)}
Koordinat $\bar{x}$ ditentukan oleh rasio antara momen statis terhadap sumbu-y ($M_y$) dengan Luas Permukaan kulit benda putar ($A$):
\[\bar{x} = \frac{M_y}{A}\]

---

\subsection*{2. Menghitung Luas Permukaan Kulit Benda Putar ($A$)}
Rumus luas permukaan benda putar terhadap sumbu-x adalah:
\[A = \int_{a}^{b} 2\pi y \, ds\]

Substitusikan $y = x^3$ dan $ds$:
\[A = 2\pi \int_{1}^{2} x^3 \sqrt{1 + 9x^4} \, dx\]

Gunakan teknik substitusi ($u$-sub):
Letakkan $u = 1 + 9x^4$, sehingga $du = 36x^3 \, dx \implies x^3 \, dx = \frac{1}{36} du$.

Perubahan batas integral:
\begin{itemize}
  \item Untuk $x = 1 \implies u = 1 + 9(1)^4 = 10$
  \item Untuk $x = 2 \implies u = 1 + 9(2)^4 = 145$
\end{itemize}

Maka nilai $A$:
\[A = 2\pi \cdot \frac{1}{36} \int_{10}^{145} \sqrt{u} \, du \]
\[A = \frac{\pi}{18} \left[ \frac{2}{3} u^{3/2} \right]_{10}^{145} \]
\[A = \frac{\pi}{27} \left( 145\sqrt{145} - 10\sqrt{10} \right) \approx 63.498\pi\]

\subsection*{3. Menghitung Momen Statis terhadap Sumbu-y ($M_y$)}
Rumus momen statis untuk kulit benda putar adalah:
\[M_y = \int_{a}^{b} 2\pi y \cdot x \, ds\]

Substitusikan komponen fungsi ke dalam rumus:
\[M_y = 2\pi \int_{1}^{2} x \cdot (x^3) \sqrt{1 + 9x^4} \, dx \]
\[M_y = 2\pi \int_{1}^{2} x^4 \sqrt{1 + 9x^4} \, dx\]

Menggunakan pendekatan integrasi numerik untuk bentuk non-elementer di atas (sebenarnya bisa diselesaikan dengan integral parsial dan substitusi fungsi trigonometri hiperbolik), diperoleh:
\[\int_{1}^{2} x^4 \sqrt{1 + 9x^4} \, dx \approx 27.661\]
\[M_y \approx 2\pi \times 27.661 = 55.322\pi\]

\subsection*{4. Menghitung Koordinat $\bar{x}$}
Substitusikan nilai $M_y$ dan $A$ ke dalam rumus titik berat:
\[\bar{x} = \frac{M_y}{A} = \frac{55.322\pi}{63.498\pi} \approx 1.162\]

\subsection*{Kesimpulan}
Titik berat dari kulit benda putar tersebut berada pada koordinat:
\[(\bar{x}, \, \bar{y}) = (1.162, \, 0)\]

\end{document}